\setcounter{Exercise}{0}

\section*{TD NSI : Choisir et manipuler les types construits}

\begin{Exercise}[title={Le bon choix de structure}]
Pour chaque situation ci-dessous, indiquez s'il est préférable d'utiliser une \textbf{liste} ou un \textbf{dictionnaire} en justifiant brièvement :
\begin{enumerate}
    \item Stocker le classement des 5 meilleurs scores d'un joueur, du premier au dernier.
    \item Stocker un carnet d'adresses où chaque nom est associé à un numéro de téléphone.
    \item Stocker les températures relevées chaque heure pendant une journée.
    \item Stocker les caractéristiques d'un personnage de jeu vidéo (nom, points de vie, inventaire).
\end{enumerate}
\end{Exercise}

\begin{Exercise}[title={Gestion d'un inventaire (Listes)}]
Un magasin de jeux vidéo reçoit un carton. En Python, on représente cet arrivage par : 
\texttt{arrivage = ["Zelda", "Mario", "Minecraft"]}.
\begin{enumerate}
    \item À l'aide d'une méthode vue en cours, ajoutez le jeu \texttt{"Tetris"} à la fin de cette liste.
    \item Le jeu \texttt{"Mario"} est abîmé. Écrivez l'instruction permettant de remplacer l'élément à l'index 1 par \texttt{"Splatoon"}.
    \item Quelle fonction permet d'afficher le nombre total d'éléments dans la liste \texttt{arrivage} ?.
\end{enumerate}
\end{Exercise}

\begin{Exercise}[title={Le Catalogue des prix (Dictionnaires)}]
Le catalogue du magasin est stocké ainsi : \\
\texttt{prix\_jeux = \{"Zelda": 60, "Minecraft": 25, "Tetris": 10\}}.
\begin{enumerate}
    \item Quelle instruction permet d'accéder au prix de \texttt{"Zelda"} ?.
    \item Une promotion est appliquée sur \texttt{"Tetris"}. Modifiez sa valeur pour qu'elle devienne \texttt{8}.
    \item Ajoutez le jeu \texttt{"Splatoon"} avec un prix de \texttt{45} au dictionnaire.
\end{enumerate}
\end{Exercise}

\begin{Exercise}[title={Algorithmes de recherche}]
On considère le dictionnaire suivant représentant le stock d'accessoires : \\
\texttt{stock = \{"Manette": 50, "Casque": 80, "Cable": 15, "Housse": 20\}}
\begin{enumerate}
    \item Écrivez un programme qui parcourt ce dictionnaire et affiche uniquement les noms des articles (clés) dont le prix est inférieur à \texttt{30}.
    \item Écrivez un programme qui calcule et affiche la somme totale des prix de tous les accessoires du dictionnaire.
\end{enumerate}
\end{Exercise}

\begin{Exercise}[title={Bonus : Structures imbriquées}]
On décide de stocker le prix et la quantité disponible dans une liste associée à chaque clé : \\
\texttt{magasin = \{"Portal": [15, 4], "Doom": [20, 0], "Roblox": [0, 100]\}}
\begin{enumerate}
    \item Quelle instruction permet d'accéder à la quantité (index 1) du jeu \texttt{"Portal"} ?
    \item Écrivez un script qui parcourt \texttt{magasin} et affiche \texttt{"RUPTURE"} si la quantité d'un jeu est égale à 0.
\end{enumerate}
\end{Exercise}